\documentclass{article}
\usepackage{amsmath, amssymb}

\title{Efficient Zero-Knowledge Proofs for Neural Network Inference}
\author{Example Authors}

\begin{document}
\maketitle

\begin{abstract}
We present a framework for generating zero-knowledge proofs of neural network inference.
Our approach handles the key operations in transformer models: matrix multiplication,
layer normalization, softmax attention, and GELU activations. We demonstrate that our
circuit construction achieves practical proving times while maintaining soundness.
\end{abstract}

\section{Introduction}

Zero-knowledge machine learning (zkML) enables a prover to demonstrate that a neural
network inference was performed correctly without revealing the model weights or input
data. A key challenge is efficiently implementing non-polynomial operations such as
Softmax and LayerNorm within arithmetic circuits.

\section{Model Architecture}

We consider a standard transformer architecture with the following components:

\begin{enumerate}
    \item Multi-head self-attention with $h$ heads and dimension $d_k = d/h$
    \item Feed-forward network with GELU activation
    \item Layer normalization after each sub-layer
    \item Residual connections
\end{enumerate}

\section{Arithmetic Circuit Construction}

\subsection{Matrix Multiplication}

For input $X \in \mathbb{R}^{n \times d}$ and weight $W \in \mathbb{R}^{d \times m}$:

\begin{equation}
    Y_{ij} = \sum_{k=1}^{d} X_{ik} \cdot W_{kj}
\end{equation}

Each multiplication-addition is implemented as a single constraint in the arithmetic circuit.
The total constraint count for MatMul is $O(n \cdot d \cdot m)$.

\subsection{Softmax Approximation}

The softmax function is defined as:

\begin{equation}
    \operatorname{Softmax}(x_i) = \frac{e^{x_i}}{\sum_{j} e^{x_j}}
\end{equation}

We approximate $e^x$ using a piecewise-linear function with $K=8$ segments over the
range $[-8, 8]$. The approximation error is bounded by $\epsilon \leq 0.01$ for inputs
in this range.

Subject to: $\sum_i \operatorname{Softmax}(x_i) = 1$

\subsection{Layer Normalization}

\begin{equation}
    \operatorname{LayerNorm}(x) = \gamma \cdot \frac{x - \mu}{\sqrt{\sigma^2 + \epsilon}} + \beta
\end{equation}

where $\mu = \frac{1}{n}\sum_{i} x_i$ and $\sigma^2 = \frac{1}{n}\sum_i (x_i - \mu)^2$.

We implement division using Newton's method with 3 iterations, and inverse square root
using a lookup table with 256 entries.

\subsection{GELU Activation}

\begin{equation}
    \operatorname{GELU}(x) \approx 0.5x\left(1 + \tanh\left[\sqrt{2/\pi}(x + 0.044715x^3)\right]\right)
\end{equation}

We use the tanh approximation form and implement tanh via a lookup table.

\subsection{ReLU Activation}

\begin{equation}
    \operatorname{ReLU}(x) = \max(0, x)
\end{equation}

Implemented using a sign decomposition: $x = x^+ - x^-$ with range checks on both components.

\section{Quantization}

All computations use 16-bit fixed-point arithmetic with 8 fractional bits. Model weights
are quantized using symmetric per-tensor quantization.

The quantization error is bounded by $2^{-8}$ per operation.

\section{Security Analysis}

\begin{theorem}
The proof system satisfies computational soundness under the discrete logarithm assumption.
The prover cannot produce a valid proof for an incorrect inference result except with
negligible probability.
\end{theorem}

\begin{theorem}
The proof system satisfies zero-knowledge: the verifier learns nothing about the model
weights or input beyond the inference result.
\end{theorem}

Constraint: all model weights must be committed using Poseidon hash before proof generation.

Constraint: every intermediate activation must be constrained to prevent layer-skipping attacks.

Constraint: range checks enforce $x \in [-2^{15}, 2^{15}-1]$ for all intermediate values.

\section{Results}

Our implementation achieves a proving time of 12 seconds for a 6-layer transformer
with $d=256$ on a single GPU, with a proof size of 1.2 KB.

\end{document}
